\documentclass{article}
\usepackage[T2A]{fontenc} % russian letters
\usepackage{algorithm}
\usepackage{algpseudocode}
\usepackage{graphicx}
\usepackage{geometry}
\geometry{
    top=1cm,
    bottom=2cm,
    left=4cm,
    right=4cm
}

\renewcommand{\figurename}{Рисунок}

\title{Лабораторная работа №1}
\author{Денис Лебедев, 2 курс, 9 группа}
\date{17 сентября 2023 г.}

\begin{document}
    \maketitle
    \section{Общие сведения}
        \hspace*{1cm}\textbf{Дихотомия} - это алгоритмический метод, который используется для нахождения решения задачи путем последовательного деления области поиска на две части и выбора той части, в которой находится решение. Этот метод широко используется в различных областях, включая математику, информатику и численные методы.
        \hspace*{1cm}\textbf{Сложность алгоритма} - количество итераций необходимых для решения задачи этим алгоритмом. В случае дихотомии область поиска решения сужается вдвое после каждой итерации, таким образом сложность алгоритма равна O($log(n)$).
        \hspace*{1cm}Ниже представлен псевдокод алгоритма нахождения загаданного числа в массиве с помощью дихотомии, а также код на языке программирования C++, содержащем элементы $[1, 2,..., num]$.

    \section{Псевдокод}
        \begin{algorithm}
        \large{\textbf{Условные обозначения: \\}}
        \small{
            $num$ - число элементов массива (массив содержит элементы $[1, 2,..., num]$)\\
            $hiddenNumber$ - загаданное число\\
            $floor$ - функция округления: возвращает наибольшее целое число, которое меньше или равно данному числу\\
            $half$ - половина элементов (число) с округлением к наименьшему целому числу \\
            $middle$ - элемент - середина интервала\\
            $p1$, $p1$ - рассматриваемый интервал $[p1, p2]$\\
            $guessNumber$ - отгаданное число\\
        } \\
        \hspace*{\algorithmicindent} \textbf{Вход: } $num$, $hiddenNumber$\\
        \hspace*{\algorithmicindent} \textbf{Выход: } $guessNumber$
        \begin{algorithmic} [1]
            \State{$p1 \gets 1$}
            \State{$p2 \gets num$}
            \While{$p1 \ne p2$}
                \State{$half \gets floor((p2 - p1 + 1) / 2)$}
                \State{$middle \gets p1 + half$}
                \If{$hiddenNumber \in [p1, middle - 1]$}
                    \State{$p2 \gets middle - 1$}
                \Else
                    \State{$p1 \gets middle$}
                \EndIf
            \EndWhile
            \State{$guessNumber \gets p1$}
            \State{\Return $guessNumber$}
        \end{algorithmic}
        \end{algorithm}

    \newpage
    \section{Код на языке C++}
        \begin{figure} [h]
            \centering
            \includegraphics[width=1\textwidth]{dihotomia.png}
            \caption{Реализация алгоритма дихотомии на языке C++}
        \end{figure}
\end{document}
